% =============================================================================
%  Probabilistic Robotics via Rust — The FCP Way
%  Print edition, in the Shannon Robotics design language.
%
%  Build:  make        (or: pdflatex main.tex, three times)
% =============================================================================
\documentclass[11pt, twoside, openright]{book}

\usepackage[
  paperwidth=7in, paperheight=10in,
  inner=0.95in, outer=0.75in, top=0.85in, bottom=0.95in,
  headsep=16pt, footskip=28pt
]{geometry}

\usepackage{shannon}

\hypersetup{
  pdftitle={Probabilistic Robotics via Rust --- The FCP Way},
  pdfauthor={Shannon Robotics},
  pdfsubject={Probabilistic robotics, state estimation, SLAM, Rust},
  colorlinks=true,
  linkcolor=shAccent,
  urlcolor=shSkyDark,
  citecolor=shAccent
}

\begin{document}

% =============================================================================
%  Front matter
% =============================================================================
\frontmatter
\pagestyle{empty}

\begin{titlepage}
\thispagestyle{empty}
\noindent
{\color{shTeal}\rule{\textwidth}{3pt}}
\vspace{2.2in}

\noindent
{\displayfont\fontsize{10}{12}\selectfont\color{shTealDark}%
 \MakeUppercase{Shannon Robotics}}

\vspace{10pt}
\noindent
{\displayfont\bfseries\fontsize{34}{39}\selectfont\color{shInk}%
 Probabilistic Robotics\\via Rust}

\vspace{12pt}
\noindent
{\displayfont\fontsize{16}{20}\selectfont\color{shInkMuted}%
 The FCP Way}

\vspace{20pt}
\noindent
{\color{shSlate300}\rule{2.4in}{0.6pt}}

\vspace{16pt}
\noindent
\begin{minipage}{4.2in}
{\fontsize{11}{15}\selectfont\color{shInkSecondary}%
 Complete mathematical foundations, interactive intuition, and working
 implementations in Rust --- for robots that must act without ever knowing
 exactly where they are.}
\end{minipage}

\vfill
\noindent
{\displayfont\fontsize{9}{12}\selectfont\color{shInkMuted}%
 Foundation \ \textperiodcentered\ \ Conceptual \ \textperiodcentered\ \ Practical}

\vspace{6pt}
\noindent
{\color{shTeal}\rule{\textwidth}{3pt}}
\end{titlepage}

% ----- Colophon --------------------------------------------------------------
\cleardoublepage
\thispagestyle{empty}
\vspace*{\fill}
\noindent
{\fontsize{9}{13}\selectfont\color{shInkSecondary}
\textbf{Probabilistic Robotics via Rust --- The FCP Way}\par
\vspace{5pt}
Print edition, generated from the interactive web edition.\par
\vspace{9pt}
This book takes Thrun, Burgard and Fox's \emph{Probabilistic Robotics} as its
spine and brings it forward: factor graphs as the inference back-end, estimation
on Lie groups, scan matching and pose-graph SLAM, visual-inertial odometry,
truncated signed distance fields, sampling-based stochastic control, online
POMDP solvers, and active SLAM. Five further texts serve as baselines throughout
--- Lynch \& Park's \emph{Modern Robotics} and Craig's \emph{Introduction to
Robotics} for the geometry, Niku for sensing hardware, Choset et al.'s
\emph{Principles of Robot Motion} for planning, and Spong, Hutchinson \&
Vidyasagar for modelling and control.\par
\vspace{9pt}
Every chapter of the web edition carries interactive simulations that run the
same algorithms the text derives. In this print edition each appears as a
described figure with its widget identifier, so the two editions can be read
side by side.\par
\vspace{9pt}
{\color{shInkMuted}Set in Space Grotesk and a Palatino-class serif.
Composed with \LaTeX.}\par}
\vspace*{\fill}

% ----- How to read this book -------------------------------------------------
\cleardoublepage
\pagestyle{fancy}
\chapter*{How to read this book}
\addcontentsline{toc}{chapter}{How to read this book}

Every chapter makes three passes over its subject, and they are meant to be read
in this order.

\textbf{Conceptual} comes first, because intuition should precede formalism. In
the web edition you meet an interactive figure before you meet an equation; here
that figure is described, and the description tells you what it would show you.

\textbf{Foundation} is the mathematics, in full. Definitions are precise,
assumptions are stated where they are made rather than buried, and derivations
are carried to the end. Long algebra sits in marked derivation blocks, so you
can take the result on faith the first time through and come back for the proof
later.

\textbf{Practical} is the Rust. Not pseudocode dressed up as code: real types,
real crates, code you could lift into a robot.

\section*{The running lab}

Two worlds recur throughout, so that every new method can be compared against
the last on identical ground. The \textbf{Hallway} is a one-dimensional corridor
with indistinguishable doors --- it exists because a belief over a 1-D world can
be drawn as a curve, which makes the Bayes filter visible in a way no
two-dimensional picture manages. The \textbf{Apartment} is a 2-D floorplan with
rooms, doorways and a long corridor, sensed by a simulated LiDAR, and
deliberately a little symmetric, because ambiguity is where probabilistic
methods earn their keep.

The robot is \shrobotname{Rusty}, a differential-drive rover with wheel encoders
and a laser scanner. Rusty is built in Chapter~4 and appears in every chapter
after it; by Chapter~26 it explores and maps an apartment it has never seen.

\section*{Colour carries meaning}

One convention runs through the prose, the equations, the figures and the code
comments alike. It is worth internalising before you start:

\shcolorkey

\noindent
Blue is always what you believed before. Orange is always what you believe after
moving. Green is always what a sensor said. Purple is always what you believe
after taking that sensor seriously. Grey, where it appears, is the truth the
robot does not have access to. The teal used for chapter furniture is
deliberately outside this set, so it can never be mistaken for data.

% ----- Contents --------------------------------------------------------------
\cleardoublepage
\setcounter{tocdepth}{0}
\renewcommand{\contentsname}{Contents}
\tableofcontents

% =============================================================================
%  Main matter
% =============================================================================
\mainmatter

\part{Foundations: The Robot and Its Uncertainty}
\input{chapters/ch01}
\input{chapters/ch02}
\input{chapters/ch03}
\input{chapters/ch04}

\part{The Bayes Filter Family}
\input{chapters/ch05}
\input{chapters/ch06}
\input{chapters/ch07}
\input{chapters/ch08}

\part{Probabilistic Models}
\input{chapters/ch09}
\input{chapters/ch10}

\part{Localization}
\input{chapters/ch11}
\input{chapters/ch12}

\part{Mapping and SLAM}
\input{chapters/ch13}
\input{chapters/ch14}
\input{chapters/ch15}
\input{chapters/ch16}
\input{chapters/ch17}
\input{chapters/ch18}
\input{chapters/ch19}

\part{Planning and Acting under Uncertainty}
\input{chapters/ch20}
\input{chapters/ch21}
\input{chapters/ch22}
\input{chapters/ch23}
\input{chapters/ch24}

\part{Frontiers and Integration}
\input{chapters/ch25}
\input{chapters/ch26}

\end{document}
